\section{A methodology for credible PNT-resilience comparison}\label{sec:method}

The difficulty in comparing resilient-PNT technologies is rarely the physics; it
is establishing that two reported numbers were produced under conditions
comparable enough to be set against each other. A claimed advantage for a
quantum sensor over a classical one may reflect the sensor, or it may reflect a
different propagator, a different error model, an unstated parameter choice, or
simply a different random seed. We therefore treat the comparison protocol
itself as the object of design, and require any quantitative claim emitted by
the engine to satisfy a three-part contract. The contract is not a code-quality
aspiration; each clause is mechanically enforced, and a result that violates any
clause is not reported.

\paragraph{(a) Reproducible from scenario, seed, and engine version.}
Every figure of merit is a deterministic function of a committed scenario
specification, an explicit pseudo-random seed, and a pinned engine version.
Given those three inputs the engine produces \emph{bit-identical} output: floating-point
results are reproduced exactly across re-runs on the same target, not merely to
within a tolerance. Determinism is a property of the implementation
(seeded generators threaded explicitly, no wall-clock or
unordered-iteration dependence in the numerical path), so a third party can
regenerate any reported value from the artifacts alone. This makes a claim
\emph{auditable}: a referee who disputes a number reruns the named scenario at
the named seed and version and obtains the same bits, or finds that they do not.
The exact commands to regenerate every figure and headline number in this paper
are listed in the reproducibility appendix.

\paragraph{(b) Every figure of merit is labelled validated or not-modeled.}
We attach to each reported quantity a machine-checkable status drawn from a
fixed vocabulary: \emph{validated} (the quantity is anchored to an external
oracle or a real reference dataset, with the anchor named), or \emph{not-modeled}
(the quantity is a model projection with no external anchor, and the absence of
an anchor is declared). The status is carried alongside the number rather than
left to prose, so the validation tier of any figure of merit is recoverable
without trusting the narrative. This discipline is what keeps the
trust evidence (Section~\ref{sec:validation}, Table~\ref{tab:validation})
honestly separated from the projections: a frames check against SOFA/ERFA and a
cold-atom-IMU dead-reckoning curve are both numbers the engine emits, but only
the first is externally anchored, and the label says so. In particular, the
quantum-sensor figures throughout this paper are performance-model projections
from cited device coefficients, not measurements; they carry the not-modeled
label and a technology-readiness tag wherever they appear
(Section~\ref{sec:crossover}).

\paragraph{(c) The quantum-versus-classical axis is a single neutral code path.}
The comparison that is the subject of Section~\ref{sec:crossover} is performed by
\emph{the same code} evaluated with different published coefficients, not by two
separate implementations. A cold-atom accelerometer and a navigation-grade IMU
are driven through one dead-reckoning path; an optical-lattice clock and a
chip-scale atomic clock through one holdover path. Only the device coefficients
differ, sourced from the respective device literature. The consequence is that
any difference in the resulting figure of merit is attributable to the sensor
parameters, not to two divergent code paths that might differ in integration
scheme, error accounting, or numerical conditioning. This removes the most
common confound in cross-technology PNT comparison, in which the quantum and
classical estimates come from different tools and the gap cannot be cleanly
ascribed.

\paragraph{Falsifiability via hand-derived tests.}
The three clauses make results reproducible and honestly labelled, but
reproducibility alone does not establish correctness: an engine can deterministically
reproduce a wrong number. We therefore require that the expected
value in every test be derived by hand from physics or from an external
reference, never read back from the engine. Frame-rotation expectations come
from closed-form rotations and the SOFA/ERFA reference values; clock-stability
expectations from analytic Allan-variance relations and NIST/Stable32 reference
data; geometric-dilution expectations from the closed-form GDOP. Because the
oracle is independent of the implementation, a regression that silently changes a
result is caught: the engine is falsifiable against an external standard rather
than self-consistent with its own prior output.

\paragraph{What is new relative to prior reproducible-research work.}
Reproducibility and model-credibility frameworks for this domain exist, and we
build on them. \citet{fernandezprades2018reproducibility} establish continuous
reproducibility for GNSS signal processing, but their scope is the raw-IF
receiver chain; the resilient-PNT figures of merit considered here (inertial
dead reckoning, clock holdover, time-transfer budgets, orbit residuals) lie
outside it. \citet{nasastd7009} provide a thorough modelling-and-simulation
credibility regime, but it is a heavyweight governance process intended for
agency programs, not a lightweight protocol a single open tool can enforce on
every emitted number. To our knowledge, no open tool combines, across the full
resilient-PNT stack, the two elements we treat as load-bearing: the
\emph{single-code-path parameter swap} as the required mechanism for any
quantum-versus-classical comparison, and the \emph{per-figure-of-merit,
machine-checkable validated/not-modeled label} carried with each result. Those
two requirements, applied together across the stack, are what is new here; the
engine of Section~\ref{sec:arch} is one instantiation of them.
